% Clase 21 — Dorantes Mejía Emmanuel (potenciómetro / ADC)
\paragraph{Potenciómetro (ADC)}
\textbf{Autor:} Dorantes Mejía Emmanuel

Se desarrollaron los archivos \texttt{read\_adc.c} y \texttt{read\_adc.h} para implementar la adquisición continua de
señales analógicas mediante el ADC de 12 bits de la Raspberry Pi Pico.

La implementación configura los canales ADC0, ADC1, ADC2 y ADC3 mediante muestreo \textit{round-robin}, con una
frecuencia configurada de 4000 muestras por segundo por canal. Se utilizan dos canales DMA encadenados para transferir
automáticamente las muestras hacia un búfer circular de cuatro elementos de 16 bits, reduciendo la intervención de la CPU.

El archivo \texttt{read\_adc.c} contiene la inicialización del ADC, la configuración del DMA y la función
\texttt{read\_adc\_is\_alive()}, que permite supervisar el avance de las transferencias para el \textit{watchdog}.
Por su parte, \texttt{read\_adc.h} define las constantes de configuración y declara las funciones del módulo.

\textbf{Driver:} se usan las bibliotecas \texttt{hardware/adc} y \texttt{hardware/dma} del Pico SDK, que escriben
directo en los registros del microcontrolador; no hay kernel ni driver de Linux en la Pico. No hace falta un driver
propio porque el ADC del RP2040 ya ofrece \textit{round-robin}, FIFO y petición de DMA (DREQ).
